\documentclass[11pt,a4paper]{article}
\usepackage[utf8]{inputenc}
\usepackage{graphicx}
\usepackage[left=2.5cm,top=3cm,right=2.5cm,bottom=3cm,bindingoffset=0.5cm]{geometry}
\usepackage{AEDLogica, AEDEspecificacion, AEDTADs}
\usepackage{caratula}

\newlength{\miSangria}
\setlength{\miSangria}{2em}

\titulo{Trabajo práctico 1: Especificacion de Tads}
\subtitulo{Tad AirMiles}

\fecha{\today}

\materia{Algoritmos y Estructuras de Datos}
\grupo{Grupo: Arraydecuatroposiciones} 

\integrante{Javier M. Ramos}{288/20}{ramosjavier276@gmail.com}
\integrante{Nancy Ramos}{1206/24}{Ramosnancydaniela@outlook.com}
\integrante{Elian Uba}{1420/23}{elianuba2004@gmail.com }
% \integrante{Apellido, Nombre2}{002/01}{email2@dominio.com} %


% Declaramos donde van a estar las figuras
% No es obligatorio, pero suele ser comodo
\graphicspath{{../static/}} 

% Asi pueden escribir nuevos comandos. 
% Este por ejemplo asegura q los nombres 
% que figuren con una tipografia diferenciada  
\newcommand{\Tipo}[1]{\mathsf{#1}}
% la sintaxis es \newcommand{\nombreDeLaMacro}[cantidadDeParametros]{Lo que va ser remplazado por el macro} 
\newcommand{\norm}[1]{\vert#1\vert}



\begin{document}

\maketitle

\section{Macros para escribir en LaTeX}

Para facilitar la escritura del trabajo práctico en LaTeX, les proveemos de varias macros que les permitirán escribir especificaciones en lógica de primer orden. 1
Por ejemplo, podrán hacer uso de comandos que definen los conectores lógicos: $\land, \lor, \implica$ y también los operadores de la lógica trivaluada $\yLuego$, $\oLuego$, $\implicaLuego$.

Además, disponen de comandos para cuantificadores: $\paraTodo{x}{\mathbb{Z}}{x \geq 0 \lor x < 0}$, $\existe{x}{\mathbb{Z}}{x \geq 0}$. A estos comandos se les puede agregar un \texttt{Largo} al final para que ocupen múltiples renglones.

\section{Renombres de Tipos y Estructuras}

FactorDePromocion \textbf{ES}: $\mathbb{Z}$

\noindent estadoPromocion \textbf{ES}: bool

\noindent topeDePromocion \textbf{ES}: $\mathbb{R}$


\noindent enum Categoría \textbf{ES} \{Bronce, Plata,Oro, Platino  \}

\vspace{0.5cm}
\noindent PropiedadesUsuario \textbf{ES} Struct \{

\hspace*{1.5em} Millas: tuple $\tupla{\begin{array}{l}
			MillasTotales: \mathbb{Z},\ MillasDisponibles: \mathbb{Z}, \\
			MillasCanjeadas: \mathbb{Z},\ MillasTransferidas: \mathbb{Z}
		\end{array}}$,

\hspace*{1.5em} VuelosRealizados: $\mathbb{Z}$,

\hspace*{1.5em} TransferenciasRealizados: $\mathbb{Z}$,

\hspace*{1.5em} Categorías: Categoría,

\hspace*{1.5em} EsReferido: bool,



\noindent \}

\vspace{0.2cm}

\noindent Transferencias \textbf{ES} Struct \{

\hspace*{1.5em} idTransferencia:  $\mathbb{Z}$,

\hspace*{1.5em} UsuarioOrigen:  $\mathbb{Z}$,

\hspace*{1.5em} UsuarioDestino:  $\mathbb{Z}$,

\hspace*{1.5em} TipoDeTransferencia:  $\mathbb{Z}$,

\hspace*{1.5em} monto:  $\mathbb{Z}$,

\noindent \}



\vspace{2em}
\begin{tad}{Airmiles}

	\obs{Usuarios}{\Diccionario{id}{PropiedadesUsuario}}
	\\
	\obs{HistorialDeTransacciones}$\seq{Transferencias}$
	\\
	\texttt{obs} Promociones : dict \{ \\
	\hspace*{1.5em} NombrePromocion: $\str$, \\
	\hspace*{1.5em} $\seq{\Tipo{FactorDePromocion}, \Tipo{estadoPromocion},\Tipo{topeDePromocion}}$ \\
	\}

	\begin{proc}{crearTAD}
		{
			\In caracteres: \seq{\cha},
			\In secuenciaDeNombres: \seq{\tupla{nombre: \seq{\cha}, apellido: \seq{\cha}}}
		}{
			\Tipo{EjemploDeTAD}
		}
		\requiere{predicadoUnilinea(caracteres)}
		\asegura{res.secuenciaDeCaracteres = caracteres}
		\aseguraLargo{ \norm{res.conjuntoDeTuplas} = \norm{secuenciaDeNombres} \land \\predicadoMultilinea(secuenciaDeNombres, res.conjuntoDeTuplas) }
	\end{proc}

	\predLargo{predicadoMultilinea}{s: \seq{\cha}, conjunto: \conj{\tupla{\seq{\cha}, \seq{\cha}}}}{
		\paraTodo{i}{\Z}{0 \leq i < \norm{s}
			\implicaLuego s[i] \in c}
	}

	\pred{predicadoUnilínea}{s: \seq{\cha}}{
		\norm{s} > 0
	}

	\aux{unAuxiliar}{k: \Z,l: \Z}{\Z}{\norm{k-l}}

\end{tad}




\end{document}